\section{Related Work}
\label{sec:related}

We organize the prior art along the two axes that distinguish the controls in this study: the \emph{security model} (negative --- enumerate and detect bad behavior --- versus positive --- enumerate and admit only authorized behavior, default-deny) and the \emph{enforcement stage} (post-generation, after the model has produced an artifact, versus at-generation, at the decoding boundary). The resulting \(2\times2\), summarized in Table~\ref{tab:related-2x2}, locates every control class and exposes the cell this work occupies. The framing is a direct port of Galloway~\cite{galloway2020}, whose true experiment contrasted a negative control (antivirus, default-allow) against a positive one (application allowlisting, default-deny) for malware execution; we ask the structurally identical question for LLM agent tool calls, scored by effect, at scale.

\subsection{Galloway as the ported template}
\label{sec:related-Host B}

Galloway~\cite{galloway2020} adopted a quantitative experimental design --- pre-test, post-test, control group --- to test whether application allowlisting prevents an effect (malware running) that signature-based antivirus, a detect-bad control, fails to prevent. The present work transplants that comparison into the agentic setting: behavioral guardrails play the role of the negative control, and grammar-constrained decoding (GCD) over a per-request authorization grammar \(G_s\) plays the role of the positive control. The dependent variable is likewise scored by effect (``did the unauthorized action occur'') rather than by any internal property of the defended system, the agentic analogue of Galloway's ``whether a piece of malware was able to run.'' Galloway could report only an empirical block rate (best 98.9\%; some malware still ran). Among the layers this study adds beyond the ported template is a formal soundness core (Section~\ref{sec:formal}): constrained generation can emit only strings in the grammar's language, and the soundness statement \(\textit{emitted} \in L(G_s) \subseteq \mathrm{Authorized}(s)\) is quantified over all logit distributions, so it holds against a fully injection-compromised model. That is what licenses a class-level elimination-by-construction claim no finite sample could --- the empirical evidence shows that behavioral controls genuinely leak at scale and that the implementation matches the theory.

\subsection{The threat: indirect prompt injection and the benchmark landscape}
\label{sec:related-threat}

The attacks we study are predominantly \emph{indirect} prompt injections delivered through a trusted data channel rather than the user message. Greshake et al.~\cite{greshake2023} characterized this class: malicious instructions ride inside content the agent retrieves and trusts (documents, emails, web pages), so the injection occupies the same input channel as any legitimate instruction. This is precisely why the in-band behavioral controls in our matrix --- the defensive system prompt (\texttt{a1}) and the embedding input classifier on the user message (\texttt{a2}) --- can be overridden: they live in the channel the attacker shares. The empirical leakiness of behavioral defenses at scale has been documented by the agent-security benchmark literature; the Agent Security Bench~\cite{zhang2024asb} evaluated a broad matrix of attack and defense methods across many models and reported high peak attack-success rates, concluding that prevention-based defenses of the detect-bad class are inadequate. Our confirmatory corpus subjects a single subject system to roughly an order of magnitude more attack trials than that benchmark; we stress that this buys tighter confidence intervals, not generalization, because the corpus is selected-on-outcome (the population of capable attacks against the undefended agent, the live-malware analogue) and bypass rates are therefore conditional on attack potency.

\subsection{Grammar-constrained and structured generation machinery}
\label{sec:related-machinery}

The mechanism we deploy --- constraining the sampler to a formal language --- builds directly on the structured-generation line of work. Willard and Louf~\cite{willard2023} cast constrained generation as an efficient finite-state-machine guided decoding problem; xgrammar~\cite{dong2024xgrammar} provides a high-throughput pushdown-automaton token-mask backend for context-free grammars; and \texttt{llama.cpp}~\cite{gerganov_llamacpp} exposes native GBNF grammars at the sampler. Our enforcement rides on this machinery, and the same C-emission-of-zero result reproduces across engines and grammar backends. We treat this cross-engine agreement as an engine-independence robustness corroboration rather than as several co-equal anchors: \texttt{vLLM}-xgrammar (the \(1.7\mathrm{B}\) blast and the steelman) and \texttt{vLLM}-guidance/llguidance (the Mistral anchor) are the confirmatory engines for the reported data, while \texttt{llama.cpp}-GBNF (on the \(35\mathrm{B}\) anchor) and SGLang-xgrammar (on the \(1.7\mathrm{B}\), old box) serve as auxiliary corroborations. The distinction we draw is one of \emph{purpose}, not of plumbing. The dominant deployed use of this machinery is structured-output \emph{shaping} --- forcing output to a JSON schema or fixed shape. Shaping is not authorization: a schema constrains the \emph{form} of a tool call, not \emph{which} URL, channel, or file the call may name. In our design the structured-output shape is held constant across every condition (it is the substrate, not a control). The novelty is per-request \emph{authorization} compiled into the grammar from the authenticated caller identity and applied at the decoder --- the action space \emph{is} the policy --- so that a forbidden parameter value is ungenerable rather than merely mis-shaped.

\subsection{The \texorpdfstring{$2\times2$}{2x2}: where the deployed controls sit, and the empty cell}
\label{sec:related-cells}

\paragraph{Negative, post-generation.}
This cell is the industry stack: defensive prompting, input classifiers, output filters, and LLM-judge / reasoning-monitor guardrails. They are probabilistic detectors of bad behavior applied after (or, for prompt-level defenses, in-band with) generation. Output filters in particular are a credential-\emph{signature} denylist, incomplete by construction: business data that carries no credential signature passes through unflagged (our \texttt{a3} arm, and the survivor class on the capable anchor in Section~\ref{sec:res-mistral}). A recurring real-world hazard in this cell is that several commercial ``agent security'' products market themselves as preventative or allowlist-based while being, mechanically, detectors --- the motivating gap this study addresses. We name these product classes in prose only; no citation is asserted for them.

\paragraph{Positive, post-generation.}
The academic frontier has moved security into this cell: deterministic, execution-time reference monitors and policy compilers that admit only authorized actions but enforce the policy \emph{after} the model has generated the tool call --- the action is produced in full, then a separate gate inspects and rejects it. Several recent systems realize this class (execution-time allowlists, policy compilers for agentic systems, and agent reference-monitor sandboxes); we describe the class in prose and do not assert citations for individual systems beyond the allowed set. Our Condition \texttt{d} is a faithful instance of exactly this class: it applies the \emph{same} allowlist that \(G_s\) compiles (\(D \equiv C\) by construction), but correctively. This is the right comparator, and it is effective at blocking the effect; we do not claim C is more secure than a perfectly wired allowlist on the blocked effect. The two controls enforce the same positive policy over the same language \(L(G_s)\) and differ in enforcement completeness, not in what they permit. The corrective allowlist re-introduces, one layer down, the failure modes positive security exists to escape: the dangerous artifact is produced (a cost and an output-channel / forensics surface --- the emission signature), a coverage obligation across every output path and consumer, and a time-of-check/time-of-use window. GCD has none of these because there is nothing to catch. Accordingly we deploy a post-parse check of this class as a defense-in-depth backstop \emph{complementing} C, not as a control C dominates.

\paragraph{Positive, at-generation --- this work.}
The remaining cell --- positive security enforced at the decoding boundary --- was empty. GCD at the sampler moves enforcement one stage earlier than the execution-time monitor: the grammar mask is applied before sampling, driving every grammar-invalid token's probability to zero, so the out-of-scope action is ungenerable by any sampler and never exists as an artifact on any surface (output buffer, logs, telemetry, stream, KV cache). The locus is the sampler, not the model: the transformer still computes a full logit distribution and under injection may assign an arbitrarily high logit to a forbidden token, but that intent never becomes a generated token. Because the grammar is compiled from authenticated identity and applied at the decoder, it is out-of-band --- nowhere the attacker's injected tokens can reach --- which is the precise reason the in-band behavioral controls fail and this one holds. This cell is where enforcement is provable and leaves no emission; with xgrammar (which removes the per-token CPU-mask cliff) C is at-par with free decoding in the agentic loop, and the per-token masking cost is reported honestly against the avoided regeneration. The enforcement locus is also distinctive: the mask is applied at every emitted token rather than after a finished call, so it acts on the free-text and streaming surfaces an execution-time monitor inspects only after the fact. This is a statement about \emph{where} enforcement happens, not a guarantee of free-text confidentiality: we scope the headline security claim to action integrity over enumerable-sink tools; free-text confidentiality leakage through channels that must remain free-text (e.g.\ \texttt{searchFiles.query}, \texttt{respondToUser.message}) is out of scope here.

\begin{table}[t]
\centering
\small
\caption{The design space as a \(2\times2\): security model \(\times\) enforcement stage. Each cell names the deployed control class; Condition labels in \texttt{typewriter}. The positive/at-generation cell is the contribution of this work.}
\label{tab:related-2x2}
\begin{tabular}{@{}p{0.14\textwidth} p{0.38\textwidth} p{0.38\textwidth}@{}}
\toprule
 & \textbf{Negative} (detect-bad) & \textbf{Positive} (allowlist, default-deny) \\
\midrule
\textbf{Post-generation} & Behavioral / LLM-judge guardrails: defensive prompt (\texttt{a1}), input classifier (\texttt{a2}), credential output filter (\texttt{a3}), full stack (\texttt{a4}); commercial ``agent security'' marketed as preventative but mechanically a probabilistic detector. & Deterministic reference monitors / policy compilers / execution-time allowlists (Condition \texttt{d} class). Same policy as \(C\); the artifact is produced then caught --- emission, coverage, and TOCTOU surfaces remain. \\
\addlinespace
\textbf{At generation} & --- & \textbf{GCD at the sampler (\texttt{c}/\texttt{c\_guided}) --- this work.} Provable (Section~\ref{sec:formal}), no emission, at-par with free decoding in the agentic loop (vLLM/xgrammar); the mask applies per token at every output surface (free-text/streaming included) rather than after a finished call. \\
\bottomrule
\end{tabular}
\end{table}

\subsection{Abliteration and the adversarial upper bound}
\label{sec:related-abliteration}

Behavioral safety can be surgically removed from a model's weights. Arditi et al.~\cite{arditi2024} showed that refusal behavior in open-weight models is mediated by a single linear direction in the residual stream that can be ablated, yielding a model that cannot refuse. This is the natural steelman for our thesis: abliteration defeats every weight- and prompt-resident behavioral defense (the negative/post-generation cell) by construction, because there is no refusal left to elicit. It cannot, however, touch a constraint that lives in the sampler. Section~\ref{sec:res-steel} reports the resulting contrast on a fully abliterated model under a top-authority malicious prompt: 97.85\% exfiltration with the sink open versus 0.00\% under a sampler-sealed grammar, with the grammar as the sole independent variable. We present this as a rebuttal exhibit rather than a primary statistic, and we are explicit (Section~\ref{sec:limits}) that the sealed cell is partly true-by-construction; its empirical content is the paired open cell (the adversary is real) and the 96.8\% provably-forced sensitive read (the engagement-evidence objection-killer), not a claim that the grammar discovered safety.
